% !TEX root = main.tex
\subsection{기존 연구와의 비교}
기능과 비용은 별도로 비교한다. 아래 문헌 비교는 원 보고서의 범위를 유지하며, 외부 구현을 이번에 재실행한 결과가 아니다.
\begin{center}\small\setlength{\tabcolsep}{6pt}
    \textbf{표 9. 관련 기능과 비교 관점}\par\smallskip
    \begin{tabular}{@{}p{2.1cm}p{6.2cm}p{5.3cm}@{}}\toprule
        \textbf{연구} & \textbf{관련 기능}         & \textbf{비교할 내용} \\ \midrule
        Sahai 등 [2] & 공급망 문서 생성·소비와 이력 추적    & 입력 참조의 공개 범위    \\
        Babel 등 [3] & 자격증명·비공개 토큰 이전·지분 소비   & DPP 생애주기와 후속 대응 \\
        GGM [4]     & 거래에 결합된 추적 자료·후속 출력 추적 & 객체 타입과 완료 조건    \\
        Azeroth [5] & 비공개 지급·거래 결합 감사 정보     & 지급과 공급망 전이의 작업량 \\
        CAP [6]     & 정책별 열람·동결·소비 후 재발행     & 권한·동결 단위와 조건    \\
        \bottomrule\end{tabular}\end{center}

표에 없는 기능을 미지원으로 판정하지 않는다. CAP의 열람은 공개 필드와 거래량 임곗값 등 정책 조건에 따른다.
기존 연구에 같은 완료 절차가 명시되지 않았다는 이유만으로, 그 절차를 추가할 수 없다고 주장하지 않는다.

Azeroth 저자 원고 Table 1(b)의 기본 설정은 MiMC7, Tree depth 32, Groth16이며 zkTransfer는 1,555,957 Gas다.
Azeroth는 잔여 금액을 ENA에 반영하지만 자체 직접 전달은 Change Note를 새로 등록한다.
증명계·Hash·작업량이 다르므로 Gas 차이를 잔여분 처리 하나의 효과로 귀속하지 않는다.
\begin{center}\small\setlength{\tabcolsep}{6pt}
    \textbf{표 10. 출처를 구분한 참고 Gas 비교}\par\smallskip
    \begin{tabular}{@{}lrl@{}}\toprule
        \textbf{대상}            & \textbf{Gas} & \textbf{범위}         \\ \midrule
        Azeroth zkTransfer     & 1,555,957    & 저자 보고값; Audit 포함    \\
        자체 직접 전달               & 2,366,228    & Audit 없음; 5회 median \\
        SharedSecret 직접 전달     & 2,724,242    & Audit 있음; 공식 1회     \\
        zkDPP Transfer         & 3,577,414    & Voucher·Change 생성   \\
        Transfer $\to$ Proceed & 5,361,496    & 두 Event median 합    \\
        \bottomrule\end{tabular}\end{center}

이 표는 동일 조건에서의 효율 순위가 아니다. Sahai의 Fabric 지표를 EVM Gas와 동일 단위로 비교하지 않는다.
원 보고서가 설명한 Babel의 운영자 원장 시간과 receipt까지 포함한 E2E도 측정 경계가 다르다.
CAP는 상태를 바꿔 재발행하는 동결을 사용하므로 현재 객체의 상태 mapping을 바꾸는 본 구현과 작업을 맞춰야 한다.

이번 결과는 구현한 Event·감사 흐름과 그 비용을 보여준다.
같은 기능·신뢰·작업량에서 외부 시스템보다 우수하다는 결론은 별도 비교가 필요하다.

